% ============================================================
%  bundle3.tex — Research Gap, Proposed Direction, Conclusion
%  % ============================================================
%  bundle3.tex — Research Gap, Proposed Direction, Conclusion
%  % ============================================================
%  bundle3.tex — Research Gap, Proposed Direction, Conclusion
%  % ============================================================
%  bundle3.tex — Research Gap, Proposed Direction, Conclusion
%  \input{sections/bundle3} from main.tex
% ============================================================

\section{Research Gap}
\label{sec:gap}

\subsection{Limitations Across Existing Studies}

The dominant fusion and detection studies evaluated across Domain~A and Domain~B operate on one or two sensor modalities simultaneously. No reviewed work addresses synchronisation fidelity when five or more heterogeneous modalities, including radar, thermal imaging, seismic sensors, acoustic arrays, and satellite imagery, contribute to a shared fusion graph at incompatible sampling rates. Multi-source fusion surveys catalogue the computational tradeoffs of data-level, feature-level, and decision-level integration~\cite{meng2023multisource}, yet treat modality count as an extensible parameter rather than a source of combinatorial complexity in temporal alignment and confidence propagation. Pairwise evaluation protocols create an illusory impression of scalability that collapses once the modality count of a deployed border system such as CIBMS is approached.

A second cross-cutting limitation concerns classification granularity and confidence calibration. Multi-agent and anomaly-detection architectures reviewed in Sections~\ref{sec:multiagent} and~\ref{sec:falsealarm} produce binary or coarse-class outputs: threat versus non-threat, or anomalous versus normal~\cite{hassan2025multiagent, agenticcyber2025}. None implements a graduated five-level classification scheme with per-level confidence calibration, formal uncertainty quantification, and documented decision boundaries separating adjacent threat tiers. This is not solely an engineering deficit; it is an operational one. Without calibrated confidence at each severity level, human operators possess no rational basis for proportional resource allocation, and response protocols collapse to a single default escalation posture regardless of actual threat magnitude.

Explainability methods surveyed in Domain~C exhibit a latency profile that is structurally incompatible with the HITL throughput requirements identified in Domain~D. SHAP and LIME generate explanation artefacts through iterative perturbation or kernel-based approximation at timescales ranging from hundreds of milliseconds to several seconds per input instance~\cite{zhang2022xai}. Military accountability frameworks stipulate that decision-makers must have access to auditable rationale before authorising escalatory action~\cite{military2024xai}, yet the alert pipelines into which these explanations must be inserted demand sub-second end-to-end delivery~\cite{hitl2023trust}. When explanation latency exceeds the operator's decision window, the HITL layer ceases to function as a deliberative safeguard and becomes instead a throughput bottleneck that delays response without improving decision quality.

\subsection{Missing Elements in Current Research}

No reviewed system models the compound failure mode in which adversarial sensor corruption at the network edge propagates through fusion, classification, explanation, and alerting stages within a single end-to-end evaluation. Physical-world adversarial attacks on object detectors are investigated in isolation on clean benchmarks~\cite{advyolo2022}, and military AI risk assessments catalogue GPS spoofing, electromagnetic jamming, and data poisoning as independent threat vectors~\cite{milai2025risk}. The distinction between digital-perturbation adversarial research, conducted on curated datasets under controlled illumination, and the physical operational reality of contested Indian border environments, where variable weather, terrain occlusion, and deliberate countermeasures co-occur, is neither acknowledged nor addressed in any integrated evaluation framework.

The surveyed literature also lacks publicly available benchmarks derived from Indian border terrain, sensor configurations, or threat typologies. Models trained and validated on standard surveillance benchmarks or NATO-oriented military imagery~\cite{satellite2024yolov8} carry embedded domain assumptions about scene geometry, illumination statistics, and threat-actor behaviour that do not transfer to Himalayan high-altitude posts, Thar Desert thermal environments, or the Sundarbans riverine terrain. This constitutes both a data-availability gap and a policy gap: operational sensor data from CIBMS deployments is necessarily classified, precluding open benchmarking, while India's defence-AI research community has not yet established synthetic or semi-synthetic equivalents calibrated to domestic operational conditions~\cite{indiaai2024springer}.

The five technical domains surveyed in Sections~\ref{sec:fusion} through~\ref{sec:distributed} are treated throughout the literature as independent research streams. No reviewed work evaluates a system architecture that jointly optimises sensor fusion fidelity, multi-agent coordination overhead, explanation latency, HITL authorisation throughput, and adversarial robustness as a coupled system of interacting constraints~\cite{meng2023multisource, flink2024kafka, uavids2025drones}. The integration gap formalised in Section~\ref{sec:introduction} is therefore confirmed as structural rather than incidental: the obstacle to unified defence surveillance is not the absence of solutions to individual sub-problems but the absence of a framework that treats their mutual interactions as first-class design constraints.


\section{Proposed Direction}
\label{sec:proposed}

The research gaps identified above motivate an architectural direction organised around four tightly coupled subsystems. The first subsystem is a hierarchical multi-modal fusion backbone in which edge inference nodes perform modality-local preprocessing and lightweight anomaly scoring at the sensor tier, transmitting compressed intermediate representations to a regional aggregation tier~\cite{flink2024kafka}. At the aggregation tier, learnable synchronisation buffers reconcile cross-modal temporal misalignments arising from disparate sampling rates across radar, thermal, seismic, acoustic, and satellite inputs. This architecture must be evaluated under five-or-more modality combinations and must quantify classification degradation under adversarial sensor dropout, extending established fusion models with explicit robustness metrics absent from current benchmarks~\cite{meng2023multisource}.

The second subsystem is a role-specialised multi-agent orchestration layer in which each agent is assigned a defined input modality, inference scope, and output schema. Agents communicate through a shared blackboard with a threat-correlation arbitrator agent responsible for producing a graduated five-level classification. The arbitrator's output must comprise per-level softmax confidence scores calibrated against a held-out validation set constructed from India-relevant operational conditions~\cite{hassan2025multiagent, agenticcyber2025}. Uncertainty quantification, whether through Monte Carlo dropout, ensemble disagreement, or conformal prediction intervals, must be a first-class output field rather than a diagnostic afterthought, enabling downstream HITL governance to reason over the arbitrator's epistemic state~\cite{zhao2022deepai}.

The third subsystem integrates explanation generation into the alert pipeline as a constrained-latency inference component. Attention-based saliency maps and compressed SHAP approximations should be generated within the same forward pass as the threat classification, bounding explanation latency to a fixed fraction of the overall inference budget~\cite{military2024xai}. These explanations must be rendered as geospatially anchored overlays on the GIS command interface~\cite{karimipour2021iot}, presenting operators with spatially contextualised rationale rather than abstract feature-attribution vectors. Interface evaluation protocols must assess operator decision quality across expertise levels to quantify and mitigate automation bias~\cite{hitl2023trust, meaningful2024hitl}.

The fourth subsystem subjects the complete pipeline, from edge sensor through stream processor to command dashboard, to a compound adversarial threat model incorporating simultaneous node failure, adversarial patch injection on visual feeds, and GPS spoofing of unmanned platforms~\cite{flink2024kafka, advyolo2022, milai2025risk}. Human authorisation must be preserved for threat classifications at Level~4 and Level~5, with interface constraints that prevent reflexive endorsement by presenting alternative classification hypotheses and associated confidence intervals to the operator~\cite{meaningful2024hitl, autonomyteaming2024}. End-to-end pipeline latency, classification accuracy, and explanation fidelity must all be reported under adversarial conditions as primary evaluation metrics, not supplementary ablations.


\section{Conclusion}
\label{sec:conclusion}

This survey has examined the defence surveillance literature across five technical domains: multi-modal sensor fusion, multi-agent AI for threat detection and classification, false-alarm reduction with explainable decision support, geospatial situational awareness with HITL governance, and real-time distributed architectures hardened against adversarial interference. Each domain contains mature component-level solutions that demonstrate strong performance on domain-specific benchmarks. The central finding, confirmed through cross-domain comparative analysis in Section~\ref{sec:comparative} and formalised in Section~\ref{sec:gap}, is that no existing architecture addresses the interactions among these five domains as coupled design constraints within a single operational framework. The integration gap is structural: fusion fidelity conditions classification confidence, explanation latency constrains HITL throughput, and adversarial corruption at the edge cascades through all downstream processing stages.

The architectural direction outlined in Section~\ref{sec:proposed} constitutes a research agenda rather than a finished system. Three open problems carry the highest priority for advancing this agenda: the construction of India-specific benchmarks calibrated to domestic terrain, sensor configurations, and threat typologies; the design and calibration of a graduated five-level threat classification scheme under realistic sensor noise and modality dropout; and the integration of bounded-latency explainability into real-time alert pipelines without degrading classification throughput. Addressing these problems in an integrated evaluation framework would represent a substantive step toward transforming India's sensor-rich but intelligence-fragmented surveillance infrastructure into a unified, graduated threat-management capability.
 from main.tex
% ============================================================

\section{Research Gap}
\label{sec:gap}

\subsection{Limitations Across Existing Studies}

The dominant fusion and detection studies evaluated across Domain~A and Domain~B operate on one or two sensor modalities simultaneously. No reviewed work addresses synchronisation fidelity when five or more heterogeneous modalities, including radar, thermal imaging, seismic sensors, acoustic arrays, and satellite imagery, contribute to a shared fusion graph at incompatible sampling rates. Multi-source fusion surveys catalogue the computational tradeoffs of data-level, feature-level, and decision-level integration~\cite{meng2023multisource}, yet treat modality count as an extensible parameter rather than a source of combinatorial complexity in temporal alignment and confidence propagation. Pairwise evaluation protocols create an illusory impression of scalability that collapses once the modality count of a deployed border system such as CIBMS is approached.

A second cross-cutting limitation concerns classification granularity and confidence calibration. Multi-agent and anomaly-detection architectures reviewed in Sections~\ref{sec:multiagent} and~\ref{sec:falsealarm} produce binary or coarse-class outputs: threat versus non-threat, or anomalous versus normal~\cite{hassan2025multiagent, agenticcyber2025}. None implements a graduated five-level classification scheme with per-level confidence calibration, formal uncertainty quantification, and documented decision boundaries separating adjacent threat tiers. This is not solely an engineering deficit; it is an operational one. Without calibrated confidence at each severity level, human operators possess no rational basis for proportional resource allocation, and response protocols collapse to a single default escalation posture regardless of actual threat magnitude.

Explainability methods surveyed in Domain~C exhibit a latency profile that is structurally incompatible with the HITL throughput requirements identified in Domain~D. SHAP and LIME generate explanation artefacts through iterative perturbation or kernel-based approximation at timescales ranging from hundreds of milliseconds to several seconds per input instance~\cite{zhang2022xai}. Military accountability frameworks stipulate that decision-makers must have access to auditable rationale before authorising escalatory action~\cite{military2024xai}, yet the alert pipelines into which these explanations must be inserted demand sub-second end-to-end delivery~\cite{hitl2023trust}. When explanation latency exceeds the operator's decision window, the HITL layer ceases to function as a deliberative safeguard and becomes instead a throughput bottleneck that delays response without improving decision quality.

\subsection{Missing Elements in Current Research}

No reviewed system models the compound failure mode in which adversarial sensor corruption at the network edge propagates through fusion, classification, explanation, and alerting stages within a single end-to-end evaluation. Physical-world adversarial attacks on object detectors are investigated in isolation on clean benchmarks~\cite{advyolo2022}, and military AI risk assessments catalogue GPS spoofing, electromagnetic jamming, and data poisoning as independent threat vectors~\cite{milai2025risk}. The distinction between digital-perturbation adversarial research, conducted on curated datasets under controlled illumination, and the physical operational reality of contested Indian border environments, where variable weather, terrain occlusion, and deliberate countermeasures co-occur, is neither acknowledged nor addressed in any integrated evaluation framework.

The surveyed literature also lacks publicly available benchmarks derived from Indian border terrain, sensor configurations, or threat typologies. Models trained and validated on standard surveillance benchmarks or NATO-oriented military imagery~\cite{satellite2024yolov8} carry embedded domain assumptions about scene geometry, illumination statistics, and threat-actor behaviour that do not transfer to Himalayan high-altitude posts, Thar Desert thermal environments, or the Sundarbans riverine terrain. This constitutes both a data-availability gap and a policy gap: operational sensor data from CIBMS deployments is necessarily classified, precluding open benchmarking, while India's defence-AI research community has not yet established synthetic or semi-synthetic equivalents calibrated to domestic operational conditions~\cite{indiaai2024springer}.

The five technical domains surveyed in Sections~\ref{sec:fusion} through~\ref{sec:distributed} are treated throughout the literature as independent research streams. No reviewed work evaluates a system architecture that jointly optimises sensor fusion fidelity, multi-agent coordination overhead, explanation latency, HITL authorisation throughput, and adversarial robustness as a coupled system of interacting constraints~\cite{meng2023multisource, flink2024kafka, uavids2025drones}. The integration gap formalised in Section~\ref{sec:introduction} is therefore confirmed as structural rather than incidental: the obstacle to unified defence surveillance is not the absence of solutions to individual sub-problems but the absence of a framework that treats their mutual interactions as first-class design constraints.


\section{Proposed Direction}
\label{sec:proposed}

The research gaps identified above motivate an architectural direction organised around four tightly coupled subsystems. The first subsystem is a hierarchical multi-modal fusion backbone in which edge inference nodes perform modality-local preprocessing and lightweight anomaly scoring at the sensor tier, transmitting compressed intermediate representations to a regional aggregation tier~\cite{flink2024kafka}. At the aggregation tier, learnable synchronisation buffers reconcile cross-modal temporal misalignments arising from disparate sampling rates across radar, thermal, seismic, acoustic, and satellite inputs. This architecture must be evaluated under five-or-more modality combinations and must quantify classification degradation under adversarial sensor dropout, extending established fusion models with explicit robustness metrics absent from current benchmarks~\cite{meng2023multisource}.

The second subsystem is a role-specialised multi-agent orchestration layer in which each agent is assigned a defined input modality, inference scope, and output schema. Agents communicate through a shared blackboard with a threat-correlation arbitrator agent responsible for producing a graduated five-level classification. The arbitrator's output must comprise per-level softmax confidence scores calibrated against a held-out validation set constructed from India-relevant operational conditions~\cite{hassan2025multiagent, agenticcyber2025}. Uncertainty quantification, whether through Monte Carlo dropout, ensemble disagreement, or conformal prediction intervals, must be a first-class output field rather than a diagnostic afterthought, enabling downstream HITL governance to reason over the arbitrator's epistemic state~\cite{zhao2022deepai}.

The third subsystem integrates explanation generation into the alert pipeline as a constrained-latency inference component. Attention-based saliency maps and compressed SHAP approximations should be generated within the same forward pass as the threat classification, bounding explanation latency to a fixed fraction of the overall inference budget~\cite{military2024xai}. These explanations must be rendered as geospatially anchored overlays on the GIS command interface~\cite{karimipour2021iot}, presenting operators with spatially contextualised rationale rather than abstract feature-attribution vectors. Interface evaluation protocols must assess operator decision quality across expertise levels to quantify and mitigate automation bias~\cite{hitl2023trust, meaningful2024hitl}.

The fourth subsystem subjects the complete pipeline, from edge sensor through stream processor to command dashboard, to a compound adversarial threat model incorporating simultaneous node failure, adversarial patch injection on visual feeds, and GPS spoofing of unmanned platforms~\cite{flink2024kafka, advyolo2022, milai2025risk}. Human authorisation must be preserved for threat classifications at Level~4 and Level~5, with interface constraints that prevent reflexive endorsement by presenting alternative classification hypotheses and associated confidence intervals to the operator~\cite{meaningful2024hitl, autonomyteaming2024}. End-to-end pipeline latency, classification accuracy, and explanation fidelity must all be reported under adversarial conditions as primary evaluation metrics, not supplementary ablations.


\section{Conclusion}
\label{sec:conclusion}

This survey has examined the defence surveillance literature across five technical domains: multi-modal sensor fusion, multi-agent AI for threat detection and classification, false-alarm reduction with explainable decision support, geospatial situational awareness with HITL governance, and real-time distributed architectures hardened against adversarial interference. Each domain contains mature component-level solutions that demonstrate strong performance on domain-specific benchmarks. The central finding, confirmed through cross-domain comparative analysis in Section~\ref{sec:comparative} and formalised in Section~\ref{sec:gap}, is that no existing architecture addresses the interactions among these five domains as coupled design constraints within a single operational framework. The integration gap is structural: fusion fidelity conditions classification confidence, explanation latency constrains HITL throughput, and adversarial corruption at the edge cascades through all downstream processing stages.

The architectural direction outlined in Section~\ref{sec:proposed} constitutes a research agenda rather than a finished system. Three open problems carry the highest priority for advancing this agenda: the construction of India-specific benchmarks calibrated to domestic terrain, sensor configurations, and threat typologies; the design and calibration of a graduated five-level threat classification scheme under realistic sensor noise and modality dropout; and the integration of bounded-latency explainability into real-time alert pipelines without degrading classification throughput. Addressing these problems in an integrated evaluation framework would represent a substantive step toward transforming India's sensor-rich but intelligence-fragmented surveillance infrastructure into a unified, graduated threat-management capability.
 from main.tex
% ============================================================

\section{Research Gap}
\label{sec:gap}

\subsection{Limitations Across Existing Studies}

The dominant fusion and detection studies evaluated across Domain~A and Domain~B operate on one or two sensor modalities simultaneously. No reviewed work addresses synchronisation fidelity when five or more heterogeneous modalities, including radar, thermal imaging, seismic sensors, acoustic arrays, and satellite imagery, contribute to a shared fusion graph at incompatible sampling rates. Multi-source fusion surveys catalogue the computational tradeoffs of data-level, feature-level, and decision-level integration~\cite{meng2023multisource}, yet treat modality count as an extensible parameter rather than a source of combinatorial complexity in temporal alignment and confidence propagation. Pairwise evaluation protocols create an illusory impression of scalability that collapses once the modality count of a deployed border system such as CIBMS is approached.

A second cross-cutting limitation concerns classification granularity and confidence calibration. Multi-agent and anomaly-detection architectures reviewed in Sections~\ref{sec:multiagent} and~\ref{sec:falsealarm} produce binary or coarse-class outputs: threat versus non-threat, or anomalous versus normal~\cite{hassan2025multiagent, agenticcyber2025}. None implements a graduated five-level classification scheme with per-level confidence calibration, formal uncertainty quantification, and documented decision boundaries separating adjacent threat tiers. This is not solely an engineering deficit; it is an operational one. Without calibrated confidence at each severity level, human operators possess no rational basis for proportional resource allocation, and response protocols collapse to a single default escalation posture regardless of actual threat magnitude.

Explainability methods surveyed in Domain~C exhibit a latency profile that is structurally incompatible with the HITL throughput requirements identified in Domain~D. SHAP and LIME generate explanation artefacts through iterative perturbation or kernel-based approximation at timescales ranging from hundreds of milliseconds to several seconds per input instance~\cite{zhang2022xai}. Military accountability frameworks stipulate that decision-makers must have access to auditable rationale before authorising escalatory action~\cite{military2024xai}, yet the alert pipelines into which these explanations must be inserted demand sub-second end-to-end delivery~\cite{hitl2023trust}. When explanation latency exceeds the operator's decision window, the HITL layer ceases to function as a deliberative safeguard and becomes instead a throughput bottleneck that delays response without improving decision quality.

\subsection{Missing Elements in Current Research}

No reviewed system models the compound failure mode in which adversarial sensor corruption at the network edge propagates through fusion, classification, explanation, and alerting stages within a single end-to-end evaluation. Physical-world adversarial attacks on object detectors are investigated in isolation on clean benchmarks~\cite{advyolo2022}, and military AI risk assessments catalogue GPS spoofing, electromagnetic jamming, and data poisoning as independent threat vectors~\cite{milai2025risk}. The distinction between digital-perturbation adversarial research, conducted on curated datasets under controlled illumination, and the physical operational reality of contested Indian border environments, where variable weather, terrain occlusion, and deliberate countermeasures co-occur, is neither acknowledged nor addressed in any integrated evaluation framework.

The surveyed literature also lacks publicly available benchmarks derived from Indian border terrain, sensor configurations, or threat typologies. Models trained and validated on standard surveillance benchmarks or NATO-oriented military imagery~\cite{satellite2024yolov8} carry embedded domain assumptions about scene geometry, illumination statistics, and threat-actor behaviour that do not transfer to Himalayan high-altitude posts, Thar Desert thermal environments, or the Sundarbans riverine terrain. This constitutes both a data-availability gap and a policy gap: operational sensor data from CIBMS deployments is necessarily classified, precluding open benchmarking, while India's defence-AI research community has not yet established synthetic or semi-synthetic equivalents calibrated to domestic operational conditions~\cite{indiaai2024springer}.

The five technical domains surveyed in Sections~\ref{sec:fusion} through~\ref{sec:distributed} are treated throughout the literature as independent research streams. No reviewed work evaluates a system architecture that jointly optimises sensor fusion fidelity, multi-agent coordination overhead, explanation latency, HITL authorisation throughput, and adversarial robustness as a coupled system of interacting constraints~\cite{meng2023multisource, flink2024kafka, uavids2025drones}. The integration gap formalised in Section~\ref{sec:introduction} is therefore confirmed as structural rather than incidental: the obstacle to unified defence surveillance is not the absence of solutions to individual sub-problems but the absence of a framework that treats their mutual interactions as first-class design constraints.


\section{Proposed Direction}
\label{sec:proposed}

The research gaps identified above motivate an architectural direction organised around four tightly coupled subsystems. The first subsystem is a hierarchical multi-modal fusion backbone in which edge inference nodes perform modality-local preprocessing and lightweight anomaly scoring at the sensor tier, transmitting compressed intermediate representations to a regional aggregation tier~\cite{flink2024kafka}. At the aggregation tier, learnable synchronisation buffers reconcile cross-modal temporal misalignments arising from disparate sampling rates across radar, thermal, seismic, acoustic, and satellite inputs. This architecture must be evaluated under five-or-more modality combinations and must quantify classification degradation under adversarial sensor dropout, extending established fusion models with explicit robustness metrics absent from current benchmarks~\cite{meng2023multisource}.

The second subsystem is a role-specialised multi-agent orchestration layer in which each agent is assigned a defined input modality, inference scope, and output schema. Agents communicate through a shared blackboard with a threat-correlation arbitrator agent responsible for producing a graduated five-level classification. The arbitrator's output must comprise per-level softmax confidence scores calibrated against a held-out validation set constructed from India-relevant operational conditions~\cite{hassan2025multiagent, agenticcyber2025}. Uncertainty quantification, whether through Monte Carlo dropout, ensemble disagreement, or conformal prediction intervals, must be a first-class output field rather than a diagnostic afterthought, enabling downstream HITL governance to reason over the arbitrator's epistemic state~\cite{zhao2022deepai}.

The third subsystem integrates explanation generation into the alert pipeline as a constrained-latency inference component. Attention-based saliency maps and compressed SHAP approximations should be generated within the same forward pass as the threat classification, bounding explanation latency to a fixed fraction of the overall inference budget~\cite{military2024xai}. These explanations must be rendered as geospatially anchored overlays on the GIS command interface~\cite{karimipour2021iot}, presenting operators with spatially contextualised rationale rather than abstract feature-attribution vectors. Interface evaluation protocols must assess operator decision quality across expertise levels to quantify and mitigate automation bias~\cite{hitl2023trust, meaningful2024hitl}.

The fourth subsystem subjects the complete pipeline, from edge sensor through stream processor to command dashboard, to a compound adversarial threat model incorporating simultaneous node failure, adversarial patch injection on visual feeds, and GPS spoofing of unmanned platforms~\cite{flink2024kafka, advyolo2022, milai2025risk}. Human authorisation must be preserved for threat classifications at Level~4 and Level~5, with interface constraints that prevent reflexive endorsement by presenting alternative classification hypotheses and associated confidence intervals to the operator~\cite{meaningful2024hitl, autonomyteaming2024}. End-to-end pipeline latency, classification accuracy, and explanation fidelity must all be reported under adversarial conditions as primary evaluation metrics, not supplementary ablations.


\section{Conclusion}
\label{sec:conclusion}

This survey has examined the defence surveillance literature across five technical domains: multi-modal sensor fusion, multi-agent AI for threat detection and classification, false-alarm reduction with explainable decision support, geospatial situational awareness with HITL governance, and real-time distributed architectures hardened against adversarial interference. Each domain contains mature component-level solutions that demonstrate strong performance on domain-specific benchmarks. The central finding, confirmed through cross-domain comparative analysis in Section~\ref{sec:comparative} and formalised in Section~\ref{sec:gap}, is that no existing architecture addresses the interactions among these five domains as coupled design constraints within a single operational framework. The integration gap is structural: fusion fidelity conditions classification confidence, explanation latency constrains HITL throughput, and adversarial corruption at the edge cascades through all downstream processing stages.

The architectural direction outlined in Section~\ref{sec:proposed} constitutes a research agenda rather than a finished system. Three open problems carry the highest priority for advancing this agenda: the construction of India-specific benchmarks calibrated to domestic terrain, sensor configurations, and threat typologies; the design and calibration of a graduated five-level threat classification scheme under realistic sensor noise and modality dropout; and the integration of bounded-latency explainability into real-time alert pipelines without degrading classification throughput. Addressing these problems in an integrated evaluation framework would represent a substantive step toward transforming India's sensor-rich but intelligence-fragmented surveillance infrastructure into a unified, graduated threat-management capability.
 from main.tex
% ============================================================

\section{Research Gap}
\label{sec:gap}

\subsection{Limitations Across Existing Studies}

The dominant fusion and detection studies evaluated across Domain~A and Domain~B operate on one or two sensor modalities simultaneously. No reviewed work addresses synchronisation fidelity when five or more heterogeneous modalities, including radar, thermal imaging, seismic sensors, acoustic arrays, and satellite imagery, contribute to a shared fusion graph at incompatible sampling rates. Multi-source fusion surveys catalogue the computational tradeoffs of data-level, feature-level, and decision-level integration~\cite{meng2023multisource}, yet treat modality count as an extensible parameter rather than a source of combinatorial complexity in temporal alignment and confidence propagation. Pairwise evaluation protocols create an illusory impression of scalability that collapses once the modality count of a deployed border system such as CIBMS is approached.

A second cross-cutting limitation concerns classification granularity and confidence calibration. Multi-agent and anomaly-detection architectures reviewed in Sections~\ref{sec:multiagent} and~\ref{sec:falsealarm} produce binary or coarse-class outputs: threat versus non-threat, or anomalous versus normal~\cite{hassan2025multiagent, agenticcyber2025}. None implements a graduated five-level classification scheme with per-level confidence calibration, formal uncertainty quantification, and documented decision boundaries separating adjacent threat tiers. This is not solely an engineering deficit; it is an operational one. Without calibrated confidence at each severity level, human operators possess no rational basis for proportional resource allocation, and response protocols collapse to a single default escalation posture regardless of actual threat magnitude.

Explainability methods surveyed in Domain~C exhibit a latency profile that is structurally incompatible with the HITL throughput requirements identified in Domain~D. SHAP and LIME generate explanation artefacts through iterative perturbation or kernel-based approximation at timescales ranging from hundreds of milliseconds to several seconds per input instance~\cite{zhang2022xai}. Military accountability frameworks stipulate that decision-makers must have access to auditable rationale before authorising escalatory action~\cite{military2024xai}, yet the alert pipelines into which these explanations must be inserted demand sub-second end-to-end delivery~\cite{hitl2023trust}. When explanation latency exceeds the operator's decision window, the HITL layer ceases to function as a deliberative safeguard and becomes instead a throughput bottleneck that delays response without improving decision quality.

\subsection{Missing Elements in Current Research}

No reviewed system models the compound failure mode in which adversarial sensor corruption at the network edge propagates through fusion, classification, explanation, and alerting stages within a single end-to-end evaluation. Physical-world adversarial attacks on object detectors are investigated in isolation on clean benchmarks~\cite{advyolo2022}, and military AI risk assessments catalogue GPS spoofing, electromagnetic jamming, and data poisoning as independent threat vectors~\cite{milai2025risk}. The distinction between digital-perturbation adversarial research, conducted on curated datasets under controlled illumination, and the physical operational reality of contested Indian border environments, where variable weather, terrain occlusion, and deliberate countermeasures co-occur, is neither acknowledged nor addressed in any integrated evaluation framework.

The surveyed literature also lacks publicly available benchmarks derived from Indian border terrain, sensor configurations, or threat typologies. Models trained and validated on standard surveillance benchmarks or NATO-oriented military imagery~\cite{satellite2024yolov8} carry embedded domain assumptions about scene geometry, illumination statistics, and threat-actor behaviour that do not transfer to Himalayan high-altitude posts, Thar Desert thermal environments, or the Sundarbans riverine terrain. This constitutes both a data-availability gap and a policy gap: operational sensor data from CIBMS deployments is necessarily classified, precluding open benchmarking, while India's defence-AI research community has not yet established synthetic or semi-synthetic equivalents calibrated to domestic operational conditions~\cite{indiaai2024springer}.

The five technical domains surveyed in Sections~\ref{sec:fusion} through~\ref{sec:distributed} are treated throughout the literature as independent research streams. No reviewed work evaluates a system architecture that jointly optimises sensor fusion fidelity, multi-agent coordination overhead, explanation latency, HITL authorisation throughput, and adversarial robustness as a coupled system of interacting constraints~\cite{meng2023multisource, flink2024kafka, uavids2025drones}. The integration gap formalised in Section~\ref{sec:introduction} is therefore confirmed as structural rather than incidental: the obstacle to unified defence surveillance is not the absence of solutions to individual sub-problems but the absence of a framework that treats their mutual interactions as first-class design constraints.


\section{Proposed Direction}
\label{sec:proposed}

The research gaps identified above motivate an architectural direction organised around four tightly coupled subsystems. The first subsystem is a hierarchical multi-modal fusion backbone in which edge inference nodes perform modality-local preprocessing and lightweight anomaly scoring at the sensor tier, transmitting compressed intermediate representations to a regional aggregation tier~\cite{flink2024kafka}. At the aggregation tier, learnable synchronisation buffers reconcile cross-modal temporal misalignments arising from disparate sampling rates across radar, thermal, seismic, acoustic, and satellite inputs. This architecture must be evaluated under five-or-more modality combinations and must quantify classification degradation under adversarial sensor dropout, extending established fusion models with explicit robustness metrics absent from current benchmarks~\cite{meng2023multisource}.

The second subsystem is a role-specialised multi-agent orchestration layer in which each agent is assigned a defined input modality, inference scope, and output schema. Agents communicate through a shared blackboard with a threat-correlation arbitrator agent responsible for producing a graduated five-level classification. The arbitrator's output must comprise per-level softmax confidence scores calibrated against a held-out validation set constructed from India-relevant operational conditions~\cite{hassan2025multiagent, agenticcyber2025}. Uncertainty quantification, whether through Monte Carlo dropout, ensemble disagreement, or conformal prediction intervals, must be a first-class output field rather than a diagnostic afterthought, enabling downstream HITL governance to reason over the arbitrator's epistemic state~\cite{zhao2022deepai}.

The third subsystem integrates explanation generation into the alert pipeline as a constrained-latency inference component. Attention-based saliency maps and compressed SHAP approximations should be generated within the same forward pass as the threat classification, bounding explanation latency to a fixed fraction of the overall inference budget~\cite{military2024xai}. These explanations must be rendered as geospatially anchored overlays on the GIS command interface~\cite{karimipour2021iot}, presenting operators with spatially contextualised rationale rather than abstract feature-attribution vectors. Interface evaluation protocols must assess operator decision quality across expertise levels to quantify and mitigate automation bias~\cite{hitl2023trust, meaningful2024hitl}.

The fourth subsystem subjects the complete pipeline, from edge sensor through stream processor to command dashboard, to a compound adversarial threat model incorporating simultaneous node failure, adversarial patch injection on visual feeds, and GPS spoofing of unmanned platforms~\cite{flink2024kafka, advyolo2022, milai2025risk}. Human authorisation must be preserved for threat classifications at Level~4 and Level~5, with interface constraints that prevent reflexive endorsement by presenting alternative classification hypotheses and associated confidence intervals to the operator~\cite{meaningful2024hitl, autonomyteaming2024}. End-to-end pipeline latency, classification accuracy, and explanation fidelity must all be reported under adversarial conditions as primary evaluation metrics, not supplementary ablations.


\section{Conclusion}
\label{sec:conclusion}

This survey has examined the defence surveillance literature across five technical domains: multi-modal sensor fusion, multi-agent AI for threat detection and classification, false-alarm reduction with explainable decision support, geospatial situational awareness with HITL governance, and real-time distributed architectures hardened against adversarial interference. Each domain contains mature component-level solutions that demonstrate strong performance on domain-specific benchmarks. The central finding, confirmed through cross-domain comparative analysis in Section~\ref{sec:comparative} and formalised in Section~\ref{sec:gap}, is that no existing architecture addresses the interactions among these five domains as coupled design constraints within a single operational framework. The integration gap is structural: fusion fidelity conditions classification confidence, explanation latency constrains HITL throughput, and adversarial corruption at the edge cascades through all downstream processing stages.

The architectural direction outlined in Section~\ref{sec:proposed} constitutes a research agenda rather than a finished system. Three open problems carry the highest priority for advancing this agenda: the construction of India-specific benchmarks calibrated to domestic terrain, sensor configurations, and threat typologies; the design and calibration of a graduated five-level threat classification scheme under realistic sensor noise and modality dropout; and the integration of bounded-latency explainability into real-time alert pipelines without degrading classification throughput. Addressing these problems in an integrated evaluation framework would represent a substantive step toward transforming India's sensor-rich but intelligence-fragmented surveillance infrastructure into a unified, graduated threat-management capability.
